% Generated draft: requires mathematical and source review.
\documentclass[11pt]{article}
\usepackage[margin=25mm]{geometry}
\usepackage{fontspec}
\setmainfont{Latin Modern Roman}
\usepackage{amsmath,amssymb,mathtools,unicode-math}
\setmathfont{Latin Modern Math}
\usepackage{enumitem,needspace,xurl,hyperref}
\setlength{\parindent}{0pt}
\setlength{\parskip}{5pt}
\setlength{\emergencystretch}{3em}
\newcommand{\N}{\mathbb{N}}
\newcommand{\Z}{\mathbb{Z}}
\newcommand{\Q}{\mathbb{Q}}
\newcommand{\R}{\mathbb{R}}
\newcommand{\C}{\mathbb{C}}
\newcommand{\F}{\mathbb{F}}
\newcommand{\A}{\mathbb{A}}
\newcommand{\PP}{\mathbb{P}}
\newcommand{\E}{\mathbb{E}}
\newcommand{\eps}{\varepsilon}
\newcommand{\dd}{\,\mathrm{d}}
\begin{document}
\section*{UnsolvedMath ID: 30005941}
\textbf{Source problem code:} catalogue entry
\subsection{Definitions and mathematical statement}
Let \(S\) and \(\beta_S\) denote a surface and a curve class in the intended scope of the supplied conjecture. The excerpt denotes its generalized Block--Göttsche polynomial by
\[
\operatorname{BPS}_{\beta_S}^{S}(q)\in\Z[q,q^{-1}],
\]
where \(q\) is a formal variable, and denotes the corresponding real Welschinger count by \(W_{\beta_S}^{S}\). Evaluation of this Laurent polynomial at \(q=-1\) is well defined. These descriptions identify the intended invariants but do not supply their complete definitions.

The conjecture attributed in the supplied excerpt to Argüz--Bousseau asks whether, for every admissible pair \((S,\beta_S)\),
\[
\operatorname{BPS}_{\beta_S}^{S}(-1)=W_{\beta_S}^{S}.
\]
Here ``admissible'' is a placeholder for hypotheses missing from the record, rather than a defined class of surfaces and curve classes. In particular, the record does not specify the real structure on \(S\), the admissible curve classes, the incidence conditions, or the conventions defining the Welschinger count. Consequently, this is a conditional transcription of the intended equality, not a fully specified universally quantified conjecture.

The background motivates these polynomials through toric Block--Göttsche Laurent polynomials, described as specializing to genus-zero complex Gromov--Witten counts at \(q=1\) and real Welschinger counts at \(q=-1\). It proposes an extension using higher-genus Gromov--Witten invariants with an insertion denoted \(\lambda_g\), for \(g\in\Z_{\geq0}\). The excerpt asserts a unique Laurent polynomial characterized by a generating-series identity after \(q=e^{iu}\), but that identity is truncated; it cannot serve here as a definition of \(\operatorname{BPS}_{\beta_S}^{S}\).

The scope of the supplied partial results is unresolved. The original excerpt asserts the equality for del Pezzo surfaces, phrased there as \(n\leq8\), without defining \(n\) in the retained passage. The supplied dated literature assessment instead reports a proof for blowups of the projective plane at at most six general points and describes an all-\(n\) conjecture as unproved. These assertions are retained as source claims; their hypotheses and relationship require review.
\subsection{Short English statement}
The conjecture attributed to Argüz--Bousseau asks whether evaluating the generalized Block--Göttsche polynomial at \(q=-1\) gives the corresponding real Welschinger count. The supplied record omits hypotheses and invariant definitions, and its accounts of the proved range differ. A precise unconditional formulation requires source review.
\subsection{Sources}
\begin{itemize}
\item Local UnsolvedMath registry, supplied complete source record: permanent ID \(30005941\), exact source problem code \texttt{catalogue entry}, title ``Welschinger Specialization of Generalized Block--Göttsche Polynomials.'' The record gives proposed year 2024 and attributes the conjecture and excerpted theorems to Argüz--Bousseau; its separate proposer field is empty.
\item \emph{Real and Logarithmic Enumerative Geometry} (2024), Oberwolfach Reports, report 23/2024, pp. 1285--1338, DOI 10.4171/owr/2024/23. Supplied report reference and literature-source link: \url{https://doi.org/10.4171/owr/2024/23}. The supplied extract contains a truncated defining identity and an asserted del Pezzo result for \(n\leq8\).
\item \emph{BPS polynomials and Welschinger invariants}, supplied literature-source link: \url{https://arxiv.org/abs/2506.02770}. The supplied assessment attributes to Arguez and Bousseau a proof of the specialization for blowups of the projective plane at at most six general points, while describing the all-\(n\) conjecture as unproved.
\item Literature assessment date: \texttt{2026-08-22}. This is supplied dated curation, not independently verified here. The record labels the problem \texttt{partially\_solved} and the statement \texttt{corrected\_verified}, and reports a check against the official Oberwolfach report. These labels and verification claims are supplied metadata, not independent mathematical verification by this draft.
\end{itemize}
\end{document}
